% !TEX root = ../main.tex
\chapter{Metodologia experimental}\label{ch:metodologia}
\fancyhead[RE]{\slshape \textsf{Capítol \thechapter.\ Metodologia experimental}}
\fancyhead[LO]{\slshape \textsf{\nouppercase \rightmark}}

\section{Principis de disseny experimental}\label{sec:principis_exp}

L'avaluació s'ha organitzat com un embut perquè cada etapa respon una pregunta
diferent. Primer es comprova que el sistema C genera artefactes vàlids sobre una
mostra ampla. Després es seleccionen casos prometedors amb proxies barats.
Només al final s'executen proves costoses: bitrate real amb CSMR i temps a
Raspberry. Aquesta estructura evita executar CSMR o hardware de baix recurs per
milers de combinacions sense evidència prèvia.

\begin{figure}[H]
\centering
\fbox{\begin{minipage}{0.92\textwidth}
\centering
compressió/descompressió base
$\rightarrow$ càlcul $\lambda$/Q
$\rightarrow$ Q-map
$\rightarrow$ full C-only
$\rightarrow$ CSMR
$\rightarrow$ Raspberry
\end{minipage}}
\caption{Seqüència metodològica, de major cobertura a major cost. El full
C-only criba candidats; CSMR mesura bitrate real; Raspberry mesura cost a bord.}
\label{fig:experiment_funnel}
\end{figure}

El compressor i el descompressor formen part del càlcul de qualitat quan es
treballa amb error mesurat: proporcionen l'error base que després es converteix
en $\lambda$/Q amb les fórmules fixed-quality. Per això els temps de compressió i
descompressió s'avaluen com una part central del sistema, no com una comprovació
accessòria. En paral·lel, els generadors de Q-map permeten repetir comparatives
àmplies de manera controlada i separar el cost de decidir la qualitat local del
cost d'aplicar el còdec neuronal.

\subsection{Terminologia experimental}

La memòria utilitza les unitats següents:

\begin{description}[leftmargin=3.1cm,style=nextline]
  \item[Retall] Imatge d'entrada de $512\times512$ píxels i vuit bandes.
  \item[Generació de Q-map] Càlcul dels 1.024 valors Q a partir d'una base de
  qualitat, una estimació d'error i, si escau, un preset.
  \item[Registre] Fila analítica d'una taula C-only, definida per retall, preset
  i política.
  \item[Subcas CSMR] Compressió real amb la referència SORTENY basada en
  TensorFlow Compression, que produeix un fitxer \code{.tfci}.
  \item[Execució temporitzada] Mesura del temps, CPU i memòria d'un binari C o
  d'un generador de Q-map.
\end{description}

També es fan servir tres conceptes de comparació:

\begin{description}[leftmargin=3.1cm,style=nextline]
  \item[Baseline] Compressió de referència contra la qual es calcula una millora
  o una pèrdua.
  \item[Màscara comuna] Mateixa ROI aplicada a reconstruccions diferents del
  mateix retall per calcular deltes justos de PSNR.
  \item[Cobertura ROI] Percentatge de blocs seleccionats pel preset. No es fixa
  manualment; és una conseqüència del detector i del llindar.
\end{description}

\section{Hipòtesis i criteris d'avaluació}\label{sec:hipotesis}

Les quatre preguntes de recerca formulades al Capítol~\ref{ch:introduccio} es
concreten en sis hipòtesis. La Taula~\ref{tab:hipotesis} fixa abans dels
resultats quina evidència s'utilitza per avaluar cadascuna.

\begin{table}[H]
\centering
\scriptsize
\begin{tabularx}{\textwidth}{p{0.6cm}p{0.5cm}X X}
\toprule
 & P & Hipòtesi & Criteri d'avaluació \\
\midrule
H1 & P1 & El pipeline C complet pot reproduir el comportament esperat del model
de referència amb memòria limitada. & Tests d'operadors, descompressió
correcta, paritat dels artefactes i comparació de reconstruccions. \\
H2 & P2 & Els presets automàtics poden substituir la selecció manual en casos
interpretables. & Cobertura espacial no degenerada, consistència entre retalls
i evidència visual de les regions seleccionades. \\
H3 & P3 & La qualitat espacial variable redueix bitrate real respecte de
\code{q204} uniforme sense perjudicar la regió protegida. & Mida dels fitxers
\code{.tfci} i PSNR de regió i fons calculats amb una mateixa màscara. \\
H4 & P3 & \code{focus\_bgq128} i \code{preserve-roi} cobreixen objectius
diferents. & Comparació taxa--distorsió entre màxim estalvi i màxima
conservació de la regió protegida. \\
H5 & P4 & Les optimitzacions C acceleren el pipeline sense alterar-ne el resultat
funcional. & Speedup en Raspberry i auditoria de paritat entre la versió de
referència C i la versió optimitzada. \\
H6 & P4 & El cost de decidir el Q-map és petit davant del cost del còdec. &
Percentatge de temps de generació del Q-map respecte de la compressió i del
pipeline complet. \\
\bottomrule
\end{tabularx}
\caption{Relació entre les preguntes de recerca, les hipòtesis i els criteris
d'avaluació.}
\label{tab:hipotesis}
\end{table}

\section{Dataset}\label{sec:dataset}

El conjunt principal conté 120 retalls procedents d'adquisicions Sentinel-2.
Cada fitxer RAW té vuit bandes, $512\times512$ píxels per banda, mostres
\code{uint16}, ordre BSQ i mida exacta de 4.194.304 bytes. Els noms dels fitxers
codifiquen tile, data d'adquisició, identificador de retall, bandes, dimensions,
datatype i endianess.

L'objectiu no és construir un benchmark global de Sentinel-2, sinó disposar de
variabilitat espacial suficient per provar cobertura, casos buits, ROI àmplies i
diferents presets. La limitació principal és espectral: no hi ha B11 ni B12. Els
presets que depenen de SWIR no poden avaluar-se en la seva versió completa.
\code{clouds} i \code{cloud\_avoid} utilitzen el fallback visible i, per tant,
els resultats no es poden llegir com una validació completa de detecció de
núvols.

\section{Baselines i polítiques}\label{sec:baselines}

\code{q204} és el baseline narratiu principal. Assigna Q=204 als 1.024 blocs,
no defineix ROI i no dona cap tractament especial a cap zona. És la referència
adequada per respondre la pregunta d'usuari: què es guanya respecte d'una
compressió uniforme? També manté continuïtat amb el punt de qualitat utilitzat
al treball previ, separant-lo de la nova escala de $\lambda$.

\code{adaptive\_s8} és un control tècnic secundari. Manté aproximadament Q
mitjana 204, però redistribueix qualitat segons dificultat relativa del bloc. No
usa cap preset. Serveix per separar l'efecte de Q variable de l'efecte de
prioritat semàntica.

\code{global\_target} es tracta a part. No és una política de ROI i, per tant,
no forma part del full factorial de presets. La seva funció metodològica és
validar la conversió global ``objectiu de PSNR/MSE $\rightarrow$ Q constant'':
si l'objectiu cau dins el rang calibrat, el sistema selecciona una Q uniforme;
si queda fora del rang mesurat, la decisió ha de saturar al límit disponible.
Això permet avaluar la precisió i els límits de la predicció abans d'entrar en
polítiques espacials.

El full C-only combina quatre polítiques analítiques. Les tres que s'interpreten
al cos de la memòria són \code{q204}, \code{adaptive\_s8} i
\code{focus\_bgq128}. La quarta és una política auxiliar de cribratge, inclosa
per explorar variants abans de fixar la ruta final; la seva definició operativa
queda documentada a l'Apèndix~\ref{app:config} i no es presenta com a mode
principal.

\section{Full C-only}\label{sec:full_c}

El full C-only és l'etapa ampla i barata de l'embut. Executa el còdec C,
registra artefactes, qualitat i proxies de compressibilitat, però no produeix
bitrate final. La seva funció és doble: comprovar robustesa sobre 120 retalls i
prioritzar quins casos mereixen CSMR.

\subsection{Disseny factorial}

El full combina:

\begin{equation}
10\;\text{presets}\times120\;\text{retalls}
\times4\;\text{polítiques}=4.800\;\text{registres}.
\end{equation}

Un \emph{registre} és una fila de resultats per una combinació de retall, preset
i política. No equival necessàriament a una compressió C nova, perquè els
controls \code{q204} i \code{adaptive\_s8} es comparteixen entre presets. Quan
una taula agregada parla de registres o d'execucions analítiques, es refereix a
aquestes files analítiques, no a 4.800 imatges diferents.

Els presets són \code{vegetation}, \code{vegetation\_green},
\code{chlorophyll}, \code{clouds}, \code{cloud\_avoid},
\code{local\_contrast}, \code{dark\_regions}, \code{low\_ndvi},
\code{high\_ndvi} i \code{water\_body}. La ruta manual queda exclosa perquè
l'objectiu del treball és substituir la selecció manual per presets automàtics.
Els llindars i la correspondència exacta amb fitxers públics es recullen a
l'Apèndix~\ref{app:config}.

Les quatre polítiques del full són \code{q204}, \code{adaptive\_s8},
\code{focus\_bgq128} i una política auxiliar de cribratge. Això produeix quatre
files analítiques per cada parella retall--preset, tot i que els dos controls
globals es reutilitzen entre presets.

\code{preserve-roi} no forma part d'aquest full factorial principal. Es va
avaluar després en una auditoria independent perquè respon una pregunta
diferent: no busca el màxim estalvi, sinó comprovar què passa quan la ROI rep una
Q alta fixa i el fons queda degradat. Aquesta separació evita barrejar el
cribratge de \code{focus\_bgq128} amb una política orientada explícitament a
conservar ROI.

\subsection{Artefactes i validacions}

El full comprova que cada Q-map té 1.024 bytes, que el bitstream C és llegible
pel descompressor, que la reconstrucció RAW recupera la mida esperada i que els
resums de Q, PSNR/MSE, entropia i zlib es poden calcular. Aquestes comprovacions
no substitueixen la validació tècnica del còdec descrita a la
Secció~\ref{sec:tests_c}; només asseguren que les combinacions experimentals
generen artefactes analitzables abans de passar a etapes més cares.

\subsection{Cribratge cap a CSMR}

Una fila candidata ha de mantenir la ROI, degradar el fons de manera mesurable i
reduir la qualitat mitjana:

\begin{align}
\Delta\mathrm{PSNR}_{ROI/adapt}&\geq-0,05\;\mathrm{dB},\\
\Delta\mathrm{PSNR}_{fons/adapt}&\leq-0,30\;\mathrm{dB},\\
\Delta\bar Q_{adapt}&<0.
\end{align}

A més, es revisen proxies de compressibilitat com entropia i zlib dels latents.
Aquests proxies no són bitrate real; només serveixen per ordenar candidats i
evitar executar CSMR sobre combinacions sense evidència prèvia. La fórmula de
ranking utilitzada per a aquesta priorització queda documentada a
l'Apèndix~\ref{app:evidence}.

\section{Auditoria de llindars}\label{sec:llindar_method}

Un llindar transforma una magnitud contínua, com NDVI o contrast visible, en una
màscara binària de ROI. La seva elecció no pot dependre només del bitrate d'un
cas: també ha d'evitar ROI buides, ROI totals, ROI massa petites i ROI massa
àmplies. L'auditoria de llindars revisa casos bons, medians i límit del full
C-only i busca un valor que mantingui cobertura interpretable en més d'una
escena.

La lectura és metodològica. Un llindar massa alt pot seleccionar només punts
aïllats i fer la ROI vulnerable al veïnat de blocs de baixa qualitat. Un llindar
massa baix pot convertir gairebé tota la imatge en ROI i eliminar l'espai on
estalviar bits. El resultat d'aquesta auditoria s'avalua al
Capítol~\ref{ch:resultats}, on es mostren els modes principals, experimentals i
\code{preserve-roi}.

\section{CSMR i bitrate real}\label{sec:csmr_method}

El bitstream C del Capítol~\ref{ch:pipeline} és un contenidor intermedi per
validar el pipeline C. Per mesurar bitrate codificat s'utilitza CSMR, la ruta de
referència de SORTENY basada en TensorFlow Compression. Aquesta ruta és la que
genera fitxers \code{.tfci}; per això és l'única etapa que permet parlar de bps
final, és a dir, bits per mostra espectral codificada.

CSMR rep els Q-maps generats en C com un array local de qualitat. L'auditoria
comprova que el tensor de qualitat empaquetat coincideix amb el Q-map sol·licitat
i després mesura bitrate, PSNR global, PSNR ROI i PSNR de fons. Aquesta etapa no
és una predicció de la taula de calibració: és una compressió real amb el
codificador de referència de SORTENY.

Per comparar polítiques es fan servir màscares comunes. Si una política
semàntica defineix una ROI, aquesta mateixa màscara s'aplica també a les
reconstruccions \code{q204} i \code{adaptive\_s8} del mateix retall. Així,
$\Delta ROI$ i $\Delta fons$ comparen exactament les mateixes mostres.

\section{Metodologia Raspberry}\label{sec:rpi_method}

La fase Raspberry mesura si la ruta C s'executa en una plataforma representativa
de baix recurs, amb quanta memòria i amb quin temps. Utilitza el mateix conjunt
de pesos, retalls representatius i llindars, però no executa TensorFlow. La
comparació s'organitza en tres blocs:

\begin{enumerate}[leftmargin=*]
  \item escalat inicial amb 1, 2 i 4 fils;
  \item versió C optimitzada amb quatre fils i auditoria de paritat;
  \item cost separat dels generadors de Q-map.
\end{enumerate}

La mètrica principal és el temps per retall i política, separat en compressió,
descompressió i total. També es registra RSS màxim, temperatura, throttling,
nombre de fils i integritat dels artefactes. Aquesta separació permet distingir
el cost del preset, que només decideix Q, del cost del còdec, que mesura error,
transforma la imatge i reconstrueix.
